\documentclass[10pt]{article}
\usepackage[T1]{fontenc}
\usepackage{lmodern,amsmath,amssymb,amsthm,microtype,booktabs}
\usepackage[margin=0.74in]{geometry}
\usepackage[colorlinks=true,urlcolor=blue,citecolor=blue]{hyperref}
\newtheorem{theorem}{Theorem}
\newtheorem{lemma}{Lemma}
\newtheorem{corollary}{Corollary}
\newcommand{\np}{\operatorname{nextprime}}
\title{Prescribed Decimal Prefixes and Prime Signatures:\\Existence and Minimum Formulas for OEIS A086558}
\author{Qichao Wang\\{\small Hebei University of Technology}\\[0.3em]
Daoyu Dong\\{\small University of Electronic Science and Technology of China}}
\date{October 26, 2026}
\linespread{0.96}
\hypersetup{pdftitle={Prescribed Decimal Prefixes and Prime Signatures},pdfauthor={Qichao Wang and Daoyu Dong}}
\begin{document}
\maketitle
\begin{abstract}
For every $n\ge1$ there is an $n$-digit integer beginning with the decimal
digits of $n$ and having the same prime signature as $n$. This proves the
nonvanishing conjecture in OEIS A086558. We give a finite formula for the
minimum using ordered prime tuples and one next-prime operation. At every
prime-power index $n=p^e$ this simplifies to the $e$-th power of the first
prime above an explicit integer root. The existence proof uses Dusart's
explicit prime-interval theorem and fifteen small prime certificates.
\end{abstract}
\small

\section{The target interval}
Let $d$ be the decimal length of $n$, and put
\begin{equation}\label{eq:LU}
 L=n10^{n-d},\qquad U=(n+1)10^{n-d}.
\end{equation}
The integers in $[L,U)$ are exactly the $n$-digit integers beginning
with $n$. A prime signature is the multiset of positive exponents in
the prime factorization; the signature of $1$ is empty.
OEIS A086558 \cite{oeis} takes the least integer in this interval with
the signature of $n$, assigning zero if the set is empty.

\begin{theorem}\label{thm:exist}
For every $n\ge1$ this set is nonempty.
\end{theorem}
\begin{proof}
The case $n=1$ has witness $1$. For $n\ge17$, let $p$ be the largest
prime factor of $n$, let $e=v_p(n)$, and put $C=n/p^e$. Define
\[
 x=(L/C)^{1/e}=p10^{(n-d)/e},\qquad
 y=(U/C)^{1/e}=x(1+1/n)^{1/e}.
\]
We will show that $x>396738$ and
\begin{equation}\label{eq:sufficient}
 25\log^2x>e(n+1).
\end{equation}
Dusart's Proposition 6.8 \cite{dusart} then supplies a prime
\[
 x<q\le x\left(1+\frac{1}{25\log^2x}\right)<y,
\]
where the last inequality follows from
\[
 (1+1/n)^{1/e}-1\ge\frac{1}{e(n+1)}.
\]
Indeed the derivative of $(1+t)^{1/e}$ on $[0,1/n]$ is at least
$1/(e(1+1/n))$, and integration proves the inequality. Since $q>x>p$,
the prime $q$ does not divide $C$. Therefore $Cq^e$ lies in $[L,U)$
and has the same signature as $n$.

For $n\ge64$, use $d\le n/4$, $e\le\log_2 n$, and the fact that
$n/\log_2n$ increases for $n>\mathrm e$. They give
\[
 x\ge10^{3n/(4\log_2n)}\ge10^8>396738.
\]
Also $(\log_2n)^3\le4n$ for $n\ge32$: the ratio decreases for
$\log n>3$ and is $125/32<4$ at $n=32$. Hence $e^3\le4n$.
Using $d\le n/2$ and $\log10>2$, we obtain
\[
 25\log^2x>\frac{25n^2}{e^2}
 \ge\frac{25}{4}ne>e(n+1),
\]
which proves \eqref{eq:sufficient}. The elementary inequality
$d\le n/4$ for $n\ge64$ follows on each interval of fixed digit length,
starting with $d=2$ and $n=64$; for $d\ge3$, $4d\le10^{d-1}$.

For $17\le n\le63$, $e\le5$. If $e=1,2,3,4,5$, respectively,
the smallest possible $n$ in this range is at least
$17,18,24,48,32$. Throughout the range $d=2$, so
$(n-d)/e\ge6$ in every case. Thus $x\ge10^6$ and
$25\log^2x>25\cdot12^2>5\cdot64\ge e(n+1)$.

The remaining cases $2\le n\le16$ have witnesses $Cq^e$ in the table
below; the fifteen displayed $q$ are primes, verified by exhaustive
trial division in the supplied program. Each displayed product has
the required signature and decimal prefix.
\begin{center}\small
\begin{tabular}{rrrr|rrrr}
\toprule
$n$&$C$&$q$&$e$&$n$&$C$&$q$&$e$\\\midrule
2&1&23&1&10&5&200000033&1\\
3&1&307&1&11&1&11000000021&1\\
4&1&67&2&12&3&200003&2\\
5&1&50021&1&13&1&1300000000021&1\\
6&3&200003&1&14&7&2000000000003&1\\
7&1&7000003&1&15&5&30000000000011&1\\
8&1&431&3&16&1&6329&4\\
9&1&30011&2&&&&\\\bottomrule
\end{tabular}
\end{center}
\end{proof}

\section{An exact finite prime-parameter formula}
For an integer $z\ge1$, let $\rho_e(z)$ be the least integer $w$ with
$w^e\ge z$, and let $\np(z)$ be the least prime strictly greater than
the integer $z$. All roots below are integer roots, including at exact
power boundaries.

\begin{theorem}[Minimum formula]\label{thm:min}
Let $n\ge2$ have exponent multiset $\{e_1,\ldots,e_t\}$. For each distinct
ordering of these exponents, and each increasing tuple of primes
$p_1<\cdots<p_{t-1}<U$, put
\[
 C=\prod_{i=1}^{t-1}p_i^{e_i},\qquad
 q=\np\left(\max\left\{p_{t-1},
             \rho_{e_t}(\lceil L/C\rceil)-1\right\}\right).
\]
When $t=1$, take $C=1$ and $p_0=1$. Form all candidates $Cq^{e_t}<U$.
Then A086558$(n)$ is their minimum. The tuple enumeration can be restricted
by the explicit integer bounds
\begin{equation}\label{eq:prune}
 \left(\prod_{i<j}p_i^{e_i}\right)p_j^{e_j+\cdots+e_t}<U
 \qquad(1\le j<t).
\end{equation}
\end{theorem}
\begin{proof}
Every candidate belongs to $[L,U)$ and has the prescribed signature,
because $q$ is distinct from the earlier primes and is the first prime
meeting both the ordering and lower-value constraints.
Conversely, factor a genuine minimum as
$p_1^{e_1}\cdots p_t^{e_t}$ with $p_1<\cdots<p_t$. Its exponent ordering
is included, its first $t-1$ primes obey \eqref{eq:prune}, and the
formula's final prime $q$ is at most $p_t$. The corresponding candidate
is no larger than the genuine minimum, so equality follows. Theorem
\ref{thm:exist} ensures a candidate exists, and every tuple range is finite.
\end{proof}

\begin{corollary}[Closed prime-power-index reduction]
If $n=p^e$ is a prime power, then
\begin{equation}\label{eq:pp}
 \mathrm{A086558}(n)=
 \left[\np\left(\rho_e(n10^{n-d})-1\right)\right]^e.
\end{equation}
For $n=1$ the value is $1$.
\end{corollary}
For example, at $n=16$, $L=16\cdot10^{14}$ and the first eligible prime
fourth root is $6329$, giving $6329^4=1604502443050081$.
Formula~\eqref{eq:pp} follows directly by ordering prime powers. The
existence theorem guarantees its candidate remains below $U$.
For $n\ge17$ this also follows from the prime interval in the proof;
the prime-power cases below $17$ are certified in the table.

The formula uses finite minimization over prime tuples; it is not a
constant-cost expression in elementary operations. Its evaluator and
independent trial-division certificates use only the Python standard
library. Initial formula outputs are compared with the published values.

\begin{thebibliography}{9}
\bibitem{oeis} A. Murthy, OEIS A086558,
\url{https://oeis.org/A086558}, accessed October 26, 2026.
\bibitem{dusart} P. Dusart, \emph{Estimates of some functions over primes
without R.H.}, arXiv:1002.0442v1 (2010), Proposition 6.8,
\url{https://arxiv.org/abs/1002.0442v1}.
\end{thebibliography}
\end{document}
